\chapter{Banco de respuestas extendido}

\section{Respuestas tecnicas}


\subsection{Why Unity Technical Artist?}


\begin{mdcode}
Bloque de codigo: text
Because my strongest work sits between realtime development, visual/asset preparation and technical communication. TwinSight required Unity interaction systems, WebGL constraints, Blender cleanup, technical visualization and user-facing inspection workflows. That intersection is closer to Unity Technical Art than to pure gameplay, pure 3D art or general web development.
\end{mdcode}


\subsection{What is your strongest project?}


\begin{mdcode}
Bloque de codigo: text
TwinSight X500. It is a Unity WebGL technical visualization prototype for drone assembly inspection. It includes component selection, exploded view, clipping tools, visual modes, information panels and usability/workload evaluation.
\end{mdcode}


\subsection{What did you personally build?}


\begin{mdcode}
Bloque de codigo: text
I worked across the prototype pipeline: preparing and organizing 3D assets, building the Unity scene and interactions, implementing UI behavior, preparing WebGL delivery, documenting the workflow and evaluating the prototype through usability/workload methods.
\end{mdcode}


\subsection{What is the technical-art part?}


\begin{mdcode}
Bloque de codigo: text
The technical-art value is the conversion of complex source geometry and documentation into an optimized, inspectable realtime experience. It includes asset cleanup, hierarchy decisions, visual modes, realtime constraints and clear presentation of technical information.
\end{mdcode}


\subsection{What would you do in production?}


\begin{mdcode}
Bloque de codigo: text
I would formalize asset naming, metadata, LODs, import validation, performance budgets, profiling, automated checks and test coverage. I would also separate prototype logic from production architecture and document pipeline standards for repeatable asset ingestion.
\end{mdcode}


\section{Respuestas de perfil}


\subsection{Are you senior?}


\begin{mdcode}
Bloque de codigo: text
I do not position myself by inflated seniority titles. I position myself by evidence. My strongest fit is Unity/realtime 3D technical visualization and technical-art pipeline work. I am looking for a role where I can contribute in that direction and grow with production experience.
\end{mdcode}


\subsection{Why not pure game development?}


\begin{mdcode}
Bloque de codigo: text
I can work with Unity and realtime interaction, but my strongest differentiator is technical visualization: turning complex technical material into interactive 3D experiences. That can exist in games, XR, simulation, product visualization or industrial tools.
\end{mdcode}


\subsection{Why should we consider you if you do not have years in industry?}


\begin{mdcode}
Bloque de codigo: text
Because the project evidence is directly relevant to roles that need Unity, realtime 3D, technical documentation, asset optimization and interactive visualization. I am transparent about what is prototype work, and I can explain how I would productionize it.
\end{mdcode}


\section{Respuestas dificiles}


\subsection{Did AI build this?}


\begin{mdcode}
Bloque de codigo: text
AI tools can support documentation, research and workflow acceleration, but I own the technical decisions, implementation direction, project framing and final outputs. I am careful not to present AI-assisted work as autonomous production AI.
\end{mdcode}


\subsection{Can you work remotely from Colombia?}


\begin{mdcode}
Bloque de codigo: text
Yes. I am based in Colombia and available for remote contractor/B2B roles. I would want to confirm timezone overlap, contract model, payment process and communication expectations.
\end{mdcode}


\subsection{Are you authorized in the EU?}


\begin{mdcode}
Bloque de codigo: text
Not currently. I expect Portuguese citizenship/passport around 2028, which may support future EU mobility, but I do not currently claim EU work authorization.
\end{mdcode}


\subsection{What compensation do you expect?}


\begin{mdcode}
Bloque de codigo: text
For aligned remote contractor work, my realistic target is around USD 2000/month. I can consider lower compensation if the role gives strong Unity/realtime 3D/technical-art experience and has a clear growth path. For higher-scope roles, I evaluate compensation based on responsibilities, risk and contract model.
\end{mdcode}


\subsection{Why accept a low offer?}


\begin{mdcode}
Bloque de codigo: text
I would only consider a lower offer if the experience is genuinely aligned with my field and helps build production credibility. I would not treat a low, unrelated role as strategically equivalent.
\end{mdcode}


\section{Preguntas para hacer al empleador}


\begin{itemize}
\item What kind of Unity projects would I work on in the first 90 days?
\item How does the team handle asset optimization and technical-art workflows?
\item Is this role more development, technical art, visualization, or content production?
\item What are the main performance constraints?
\item What does success look like after 3 months?
\item Is the contract employee, EOR, contractor or freelance?
\item Are there country or work-authorization limitations?
\item How are test assignments evaluated?
\end{itemize}
